%! Dimensions of the Four Subspaces — Unit 3, Lesson 3.5 — Slides (Beamer)
\documentclass[aspectratio=169, 11pt]{beamer}
\usepackage{linalg-beamer}   % provides \burgundyheader{} and \sectionlabel[color]{}
\newcommand{\vv}{\mathbf{v}}
\newcommand{\ww}{\mathbf{w}}
\newcommand{\xx}{\mathbf{x}}
\newcommand{\yy}{\mathbf{y}}
\newcommand{\cc}{\mathbf{c}}
\newcommand{\svec}{\mathbf{s}}
\newcommand{\zero}{\mathbf{0}}
\newcommand{\T}{\mathsf{T}}

\begin{document}

% TITLE SLIDE (CourseName is not defined in beamer, so the name is literal)
{
\setbeamercolor{background canvas}{bg=burgundy}
\begin{frame}[plain]
  \vspace{2.8cm}\hspace{0.55cm}%
  \begin{minipage}{0.9\textwidth}
    {\color{goldacc}\bfseries\small LESSON 3.5}\\[0.5cm]
    {\color{white}\bfseries\LARGE Dimensions of the Four Subspaces}\\[1.0cm]
    {\color{blushmid}\small Linear Algebra~$\cdot$~Unit 3: The Four Fundamental Subspaces}
  \end{minipage}
\end{frame}
}

% HOOK
\begin{frame}[t, plain]
  \burgundyheader{How does a matrix carve up each room?}
  \sectionlabel{Hook}
  \begin{columns}[T]
    \begin{column}{0.58\textwidth}
      A lever setting is an input $\xx\in\mathbb{R}^n$; the tip lands at $A\xx\in\mathbb{R}^m$.\\[6pt]
      So $A$ connects \textbf{two rooms}: an \emph{input room} $\mathbb{R}^n$ and an \emph{output room}
      $\mathbb{R}^m$.\\[6pt]
      Each room splits in two --- ``directions that do something'' and ``directions that vanish.'' That is
      the \textbf{four fundamental subspaces}.
    \end{column}
    \begin{column}{0.40\textwidth}
      \centering
      \begin{tikzpicture}[scale=0.62]
        \draw[rounded corners, burgundy, very thick] (-1.1,-1.1) rectangle (1.1,1.1);
        \node[burgundy] at (0,1.45) {\scriptsize input $\mathbb{R}^n$};
        \draw[rounded corners, royalblue, very thick] (2.9,-1.1) rectangle (5.1,1.1);
        \node[royalblue] at (4,1.45) {\scriptsize output $\mathbb{R}^m$};
        \draw[->, ultra thick, charcoal] (1.3,0)--(2.7,0) node[midway, above]{\scriptsize $A$};
        \node[burgundy] at (0,0.35){\scriptsize row space};
        \node[burgundy] at (0,-0.35){\scriptsize nullspace};
        \node[royalblue] at (4,0.35){\scriptsize column sp.};
        \node[royalblue] at (4,-0.35){\scriptsize left null};
      \end{tikzpicture}
    \end{column}
  \end{columns}
  \vspace{2pt}
  \begin{block}{Today}
    Count the dimension of all four subspaces --- from the single number $r$.
  \end{block}
\end{frame}

% BIG IDEA: TWO ROOMS, FOUR SUBSPACES
\begin{frame}[t, plain]
  \burgundyheader{Two rooms, four subspaces}
  \sectionlabel{The big idea}
  For an $m\times n$ matrix $A$ (rank $r$):
  \begin{itemize}
    \setlength{\itemsep}{5pt}
    \item \textbf{Input room $\mathbb{R}^n$:} the \emph{row space} $C(A^{\T})$ (combinations of the rows) and the \emph{nullspace} $N(A)$ (do-nothing inputs).
    \item \textbf{Output room $\mathbb{R}^m$:} the \emph{column space} $C(A)$ (reachable outputs) and the \emph{left nullspace} $N(A^{\T})$ (the $\yy$ with $A^{\T}\yy=\zero$).
  \end{itemize}
  \vspace{2pt}
  \begin{block}{Just transpose}
    The row space is ``$C$ of $A^{\T}$'' and the left nullspace is ``$N$ of $A^{\T}$'' --- old tools, applied to $A^{\T}$.
  \end{block}
\end{frame}

% DIMENSIONS FROM r
\begin{frame}[t, plain]
  \burgundyheader{Every dimension comes from $r$}
  \sectionlabel[goldacc]{Count}
  \[
    \dim C(A)=r, \quad \dim C(A^{\T})=r, \quad \dim N(A)=n-r, \quad \dim N(A^{\T})=m-r.
  \]
  \begin{itemize}
    \setlength{\itemsep}{4pt}
    \item \textbf{Row rank $=$ column rank:} independent rows and independent columns are the \emph{same} number $r$.
    \item Two counting rules, one per room:
      \[ r+(n-r)=n \quad\text{(input)}, \qquad r+(m-r)=m \quad\text{(output)}. \]
  \end{itemize}
\end{frame}

% ALL FOUR FROM ONE REDUCTION
\begin{frame}[t, plain]
  \burgundyheader{All four bases from one reduction}
  \sectionlabel{Read it off}
  \[
    A=\begin{bmatrix} 1 & 1 & 2 \\ 2 & 1 & 3 \\ 3 & 2 & 5 \end{bmatrix}
    \longrightarrow
    \begin{bmatrix} 1 & 0 & 1 \\ 0 & 1 & 1 \\ 0 & 0 & 0 \end{bmatrix}, \qquad r=2.
  \]
  \begin{itemize}
    \setlength{\itemsep}{3pt}
    \item \textbf{$C(A)$:} pivot columns $\begin{bsmallmatrix} 1\\2\\3 \end{bsmallmatrix},\begin{bsmallmatrix} 1\\1\\2 \end{bsmallmatrix}$ \; \textbf{$C(A^{\T})$:} rows $(1,0,1),(0,1,1)$ \quad ($\dim 2$ each).
    \item \textbf{$N(A)$:} $\svec=\begin{bsmallmatrix} -1\\-1\\1 \end{bsmallmatrix}$ \; \textbf{$N(A^{\T})$:} $\yy=\begin{bsmallmatrix} 1\\1\\-1 \end{bsmallmatrix}$ \quad ($\dim 1$ each).
    \item Checks: $2+1=n=3$ and $2+1=m=3$.
  \end{itemize}
\end{frame}

% TWO REDUNDANCIES
\begin{frame}[t, plain]
  \burgundyheader{Two nullspaces, two redundancies}
  \sectionlabel[goldacc]{Meaning}
  \begin{columns}[T]
    \begin{column}{0.5\textwidth}
      \textbf{Nullspace} $\to$ column redundancy\\[4pt]
      $\svec=\begin{bsmallmatrix} -1\\-1\\1 \end{bsmallmatrix}$ says $A\svec=\zero$:\\[3pt]
      $\mathbf{a}_3=\mathbf{a}_1+\mathbf{a}_2$ (column~3 redundant).
    \end{column}
    \begin{column}{0.5\textwidth}
      \textbf{Left nullspace} $\to$ row redundancy\\[4pt]
      $\yy=\begin{bsmallmatrix} 1\\1\\-1 \end{bsmallmatrix}$ says $A^{\T}\yy=\zero$:\\[3pt]
      row$_3=$ row$_1+$ row$_2$ (row~3 redundant).
    \end{column}
  \end{columns}
  \vspace{6pt}
  \begin{block}{Same tool, two directions}
    The nullspace names dependencies among \emph{columns}; the left nullspace names dependencies among \emph{rows}.
  \end{block}
\end{frame}

% ROADMAP
\begin{frame}[t, plain]
  \burgundyheader{Where this goes}
  \sectionlabel{Connections}
  \textbf{Today:} four subspaces, two rooms, all dimensions from $r$ --- $r,\,r,\,n-r,\,m-r$ --- read off one
  reduction.\\[8pt]
  \begin{itemize}
    \setlength{\itemsep}{4pt}
    \item \textbf{3.6} The Fundamental Theorem of Linear Algebra: assemble the four subspaces into the ``big picture,'' with the two pairs at right angles.
    \item \textbf{Unit 4} Orthogonality: the row space and nullspace really are perpendicular.
  \end{itemize}
  \vspace{4pt}
  \begin{block}{Keep in mind}
    One number $r$ controls it all: rank, row rank, column rank, and every dimension.
  \end{block}
\end{frame}

\end{document}
